\documentclass[10pt,a4paper]{amsart}
\usepackage[margin=19mm]{geometry}
\usepackage{amsmath,amssymb,mathtools}
\usepackage[scr=rsfs,cal=euler]{mathalfa}
\usepackage{xfrac}
\usepackage[unicode,hidelinks,bookmarks=false]{hyperref}
\newcommand{\ie}{i.e.\ }
\newcommand{\cf}{cf.\ }
\newcommand{\resp}{respectively\ }
\newcommand{\Subsection}{\subsection}
\providecommand{\nobreakdash}{\nobreak\hbox{-}}
\setlength{\emergencystretch}{2em}
\multlinegap0pt
\usepackage{bm}
\usepackage{float}
\newtheorem{lemma}{Lemma}
\title{Multiscaling in Wasserstein spaces}
\author{Wael Mattar, Nir Sharon}
\date{Source: arXiv:2509.10415v2 (CC BY 4.0)}
\begin{document}
\maketitle

\noindent\emph{Source and licence.} Multiscaling in Wasserstein spaces. Version arXiv:2509.10415v2 (CC BY 4.0), distributed by arXiv under the Creative Commons Attribution 4.0 International licence, http://creativecommons.org/licenses/by/4.0/, as recorded in the arXiv licence metadata for this exact version and shown on the abstract page https://arxiv.org/abs/2509.10415v2. Full licence notice: \texttt{licenses/arxiv-2509.10415-CC-BY-4.0.txt}.\par\smallskip

\noindent\textit{Editorial scope.} Counted unit: the article's lemma lem:details\_delta\_bound (source0.tex lines 417--422). The article's complete original proof of it is delivered with it (source0.tex lines 424--436, one \texttt{proof} environment of 13 lines). No mathematics is written, repaired or re-derived: the statement and the proof are the article's own lines, in the article's own notation, and no retained line is substituted or re-flowed.  Retained uncounted as the source-local context the proof needs: lines 121--123; lines 207--209; lines 278--280.  Nothing was omitted from any retained span, so the package carries no omission record.  No retained line contains a citation, so the wrapper carries no bibliography and no bibliography entry.  No retained line contains a figure environment, an image-inclusion command or drawing code, so no image is delivered and the package declares no asset dependency.  Editorial formatting only, in addition: an AMS \texttt{amsart} preamble that restates the article's own declarations and macros used by the retained lines, namely the environments \texttt{lemma} on separate counters, together with the standard packages those lines need.  The retained environments print in this document as Lemma 1 in order of appearance, whereas the article prints the same items with the numbering of its own full document.  The complete native source of the article as distributed in the arXiv e-print for this exact version is bundled unchanged as \texttt{sources/184974/source0.tex}.
\begin{equation}~\label{eqn:constant_speed}
    W_p(\mu_t, \mu_s) = (t-s)W_p(\mu_0, \mu_1), \quad 0\leq s\leq t\leq 1.
\end{equation}

\begin{equation}~\label{eqn:useful_relation}
    \mathcal{J}_p\big((I,T_{\mu}^\nu)_{\#}\mu\big) = \int_{\mathbb{R}^d}\|T_\mu^\nu(x) - x\|^p d\mu(x) = \|\nu \ominus \mu\|^p_{L^p(\mu)} = W_p^p(\mu, \nu),
\end{equation}

\begin{equation}~\label{eqn:elementary_multiscaling}
    \boldsymbol{\mu}^{(\ell-1)}=\mathcal{D}\boldsymbol{\mu}^{(\ell)}, \quad \boldsymbol{\psi}^{(\ell)}= \boldsymbol{\mu}^{(\ell)} \ominus \mathcal{S}\boldsymbol{\mu}^{(\ell-1)}, \quad \ell=1,\dots, J,
\end{equation}

\begin{lemma}~\label{lem:details_delta_bound}
     Let $\boldsymbol{\mu}^{(J)}$ be a sequence in $\mathcal{P}_p(\mathbb{R}^d)$ associated with the grid $2^{-J}\mathbb{Z}$. Then the detail coefficients $\boldsymbol{\psi}^{(\ell)}$ generated by the elementary multiscale transform~\eqref{eqn:elementary_multiscaling} satisfy
     \begin{equation}~\label{eqn:details_delta_bound}
         \|\boldsymbol{\psi}^{(\ell)}\|_{\infty} \leq 2\Delta \boldsymbol{\mu}^{(\ell)}, \quad \ell=1,\dots,J.
     \end{equation}
\end{lemma}

\begin{proof}
    The elements $\psi^{(\ell)}_{2i}$ are equal to the trivial zero map for all $\ell=1,\dots,J$ and $i\in\mathbb{Z}$. Therefore, $\|\psi^{(\ell)}_{2i}\|_{L^p(\mu^{(\ell)}_{2i})}=0$. Direct calculations of a general term $\psi^{(\ell)}_{2i+1}$ associated with an odd index give
    \begin{align*}
        \psi^{(\ell)}_{2i+1} & = \mu^{(\ell)}_{2i+1} \ominus \mathfrak{M}\big(\mu^{(\ell-1)}_{i}, \mu^{(\ell-1)}_{i+1}; \frac{1}{2}\big) = \mu^{(\ell)}_{2i+1} \ominus \mathfrak{M}\big(\mu^{(\ell)}_{2i}, \mu^{(\ell)}_{2i+2}; \frac{1}{2}\big).
    \end{align*}
    Consequently, by~\eqref{eqn:useful_relation} we get
    \begin{align*}
        \|\psi^{(\ell)}_{2i+1}\|_{L^p(\mu^{(\ell)}_{2i+1})} & = W_p\big(\mu^{(\ell)}_{2i+1},\; \mathfrak{M}(\mu^{(\ell)}_{2i}, \mu^{(\ell)}_{2i+2}; \frac{1}{2})\big) \\
        & \leq W_p\big(\mu^{(\ell)}_{2i+1}, \mu^{(\ell)}_{2i}\big) + W_p\big(\mu^{(\ell)}_{2i},\;\mathfrak{M}(\mu^{(\ell)}_{2i}, \mu^{(\ell)}_{2i+2}; \frac{1}{2})\big) \\
        & \leq \Delta \boldsymbol{\mu}^{(\ell)} + \frac{1}{2} W_p(\mu^{(\ell)}_{2i}, \mu^{(\ell)}_{2i+2}) \leq 2 \Delta \boldsymbol{\mu}^{(\ell)}.
    \end{align*}
    The first inequality is due to the metric property of $W_p$, and the second inequality is due to the constant-speed property~\eqref{eqn:constant_speed}. Taking the supremum norm over $i\in\mathbb{Z}$ gives the required result.
\end{proof}

\end{document}
